% Generado por sistema/backend ReporteTex.java -- no editar a mano.
% Datos: Base de datos: 148 importado (titulados) / Base de datos: 53 importado (empleadores)
\subsubsection*{F8. Medio de convocatoria de las organizaciones}

Verificar si las organizaciones convocan profesionales por todos los medios por igual o si hay medios preferidos.

\textit{Paso 1. Contar con la muestra y determinar su tamaño.} Población: organizaciones encuestadas; tamaño $n = 53$; opciones de la pregunta $k = 6$; $\alpha = 0.05$.

\begin{itemize}[nosep]
  \item $H_0$: Las respuestas de «Medio de convocatoria» se reparten por igual entre las 6 opciones (uniformidad).
  \item $H_1$: Las respuestas de «Medio de convocatoria» no se reparten por igual: hay opciones preferidas.
\end{itemize}

\textit{Paso 2. Calcular la frecuencia esperada.} Si las respuestas se repartieran por igual, cada opción tendría:
\[
FE = \frac{n}{k} = \frac{53}{6} = 8.83
\]

\textit{Paso 3. Calcular el estadístico chi-cuadrado.} Se cuenta la frecuencia observada $FO_i$ de cada opción y se calcula $\frac{(FO_i - FE)^2}{FE}$.

\begin{mitabla}{F8: prueba de las frecuencias, Medio de convocatoria}
\begin{tabular}{L{6cm}ccc}
\toprule
Opción & $FO_i$ & $FE$ & $\frac{(FO_i-FE)^2}{FE}$ \\
\midrule
Medios escritos/audiovisuales & 2 & 8.83 & $\frac{(2-8.83)^2}{8.83} = 5.2862$ \\[3pt]
Internet y redes sociales & 28 & 8.83 & $\frac{(28-8.83)^2}{8.83} = 41.5881$ \\[3pt]
Postulación competitiva & 8 & 8.83 & $\frac{(8-8.83)^2}{8.83} = 0.0786$ \\[3pt]
Oferta profesional personal & 2 & 8.83 & $\frac{(2-8.83)^2}{8.83} = 5.2862$ \\[3pt]
Recomendación de terceros & 5 & 8.83 & $\frac{(5-8.83)^2}{8.83} = 1.6635$ \\[3pt]
Otros & 8 & 8.83 & $\frac{(8-8.83)^2}{8.83} = 0.0786$ \\[3pt]
\midrule
Total & 53 & 53 & $x_0^2 = 53.9811$ \\
\bottomrule
\end{tabular}
\end{mitabla}

\[
x_0^2 = \sum_{i=1}^{6} \frac{(FO_i - FE)^2}{FE} = 53.9811
\]

\textit{Paso 4. Determinar el estadístico teórico.} $x^2_{\alpha,\,k-1} = x^2_{0.05,\,6-1} = x^2_{0.05,\,5} = 11.070$ (tabla chi-cuadrada).

\textit{Paso 5. Establecer la relación entre estadísticos.} $x_0^2 < x^2_{\alpha,\,k-1}$ $\rightarrow$ $53.9811 \ge 11.070$ $\rightarrow$ no.

\textit{Paso 6. Conclusión.}

\begin{center}
\begin{tabular}{l}
Si ($x_0^2 < x^2_{\alpha,\,k-1}$)\\
\quad ``Uniformidad''\\
Sino\\
\quad \fbox{No ``Uniformidad''}\\
FinSi
\end{tabular}
\end{center}

Se rechaza $H_0$: las respuestas no se reparten por igual, hay opciones preferidas. La más elegida es «Internet y redes sociales» (28 de 53, 52.8\%).

\textit{Recomendación.} Hay medios preferidos: la Carrera debe publicar sus titulados y su bolsa de trabajo en esos medios.

